\documentclass[11pt,letterpaper]{amsart}
\usepackage[letterpaper,margin=1in]{geometry}
\usepackage{amsmath,amssymb,lmodern}
\usepackage[hidelinks,unicode]{hyperref}
\usepackage{fancyvrb}
\hypersetup{pdftitle={A prime primitive root below every odd prime},
  pdfsubject={Erdős problem 985},pdfauthor={Will Howe},
  pdfkeywords={prime primitive root, Dirichlet character, Erdős problem 985}}
\newtheorem{theorem}{Theorem}[section]
\newtheorem{lemma}[theorem]{Lemma}
\newcommand{\ord}{\operatorname{ord}}
\newcommand{\gr}{\widehat g}
\title[A prime primitive root below every odd prime]{A prime primitive root below every odd prime}
\author[W. Howe]{Will Howe}
\date{November 20, 2026}
\subjclass[2020]{11A07, 11M06}
\keywords{Prime primitive root, multiplicative order, Dirichlet character}
\thanks{The author used OpenAI Codex to prepare the manuscript and verify the finite calculations.}
\begin{document}
\begin{abstract}
Every odd prime \(p\) has a prime primitive root \(q<p\). We deduce the
large-prime case from McGown, Trevi\~no, and Trudgian's bound and the
Dirichlet generalized Riemann hypothesis proved by Alp\"oge with Claude.
For the 405 odd primes at most \(2791\), we give an exhaustive calculation
and verify the multiplicative orders by a second method.
\end{abstract}
\maketitle
\section{Introduction}
Erd\H{o}s asked in 1965 whether each prime \(p\) admits a prime primitive
root strictly below \(p\) \cite[p.~233]{erdos}. The prime \(2\) is exceptional
because there is no prime below it. For \(a\) coprime to \(p\), write \(\ord_p(a)\) for
the least positive integer \(n\) such that \(a^n\equiv1\pmod p\).

\begin{theorem}[Erd\H{o}s problem 985]\label{t:main}
For every odd prime \(p\), there is a prime integer \(q\) such that
\[
                         q<p,\qquad \ord_p(q)=p-1 .
\]
\end{theorem}
An integer of order \(p-1\) generates \((\mathbb Z/p\mathbb Z)^\times\).
Let \(\gr(p)\) be the least prime integer with this property. Under GRH,
McGown, Trevi\~no, and Trudgian \cite[Theorem~1]{mtt} proved
\begin{equation}\label{eq:mtt}
                         \gr(p)<\sqrt{p}-2\qquad(p>2791,\ p\text{ prime}).
\end{equation}
Their explicit formula uses the Riemann zeta function and the Dirichlet
\(L\)-functions for characters modulo \(p\) \cite[\S2, Lemma~1]{mtt}.

Alp\"oge's formal proof, developed with Claude \cite{alp}, establishes RH
for zeta. For every positive modulus \(N\), complex Dirichlet character
\(\chi\bmod N\), and \(s\in\mathbb C\), it also gives
\[
            0<\Re s<1,\quad L(s,\chi)=0
            \quad\Longrightarrow\quad \Re s=\tfrac12 .
\]
These statements supply the hypotheses of \eqref{eq:mtt}. Sutherland reports
that all \(175{,}187\) modules passed Lean's kernel, that two independent kernels
accepted the full export, and that the exported theorems match the stated
zeta and Dirichlet propositions \cite{mit}.

\begin{proof}[Proof of Theorem~\ref{t:main}]
For \(p>2791\), choose \(q=\gr(p)\). Equation~\eqref{eq:mtt} gives
\(q<\sqrt p-2<p\), and the definition of \(\gr(p)\) gives \(\ord_p(q)=p-1\).
For \(p\le2791\), apply Lemma~\ref{l:finite} below.
\end{proof}

The conditional reduction and finite check appeared in January 2026
\cite{jan}. An earlier prize submission recorded the deduction after
Alp\"oge's announcement \cite{prior}. The accompanying
\href{https://github.com/H0W4RD-dev/erdos-985-prime-primitive-root/blob/main/endpoint.md}{calculation at \(x=p\)}
also derives the conclusion directly from the character-sum estimates
of \cite[Lemma~1]{mtt}. Its finite remainder has \(74{,}187\) candidates,
including \(18{,}802\) prime moduli.

\newpage
\section{The finite endpoint}
\begin{lemma}[Small prime moduli]\label{l:finite}
For every odd prime \(p\le2791\) a prime \(q\le53\) satisfies \(q<p\) and
\(\ord_p(q)=p-1\).
\end{lemma}
\begin{proof}
Trial division identifies the primes. Starting at \(1\), the function
\texttt{order} finds the first return to \(1\) under multiplication by \(q\)
modulo \(p\). Fermat's theorem guarantees a return by exponent \(p-1\).
\begin{Verbatim}[fontsize=\small,xleftmargin=1.7em]
from math import isqrt

def prime(n):
    return n >= 2 and all(n % d for d in range(2, isqrt(n) + 1))

def order(q, p):
    a = 1
    for k in range(1, p):
        a = (a * q) % p
        if a == 1:
            return k

bases = [q for q in range(2, 54) if prime(q)]
roots = []
for p in range(3, 2792, 2):
    if prime(p):
        q = next(q for q in bases if q < p and order(q, p) == p-1)
        roots.append((p, q))
assert (len(roots), max(q for p, q in roots)) == (405, 53)
\end{Verbatim}
The calculation produces all 405 witnesses. A second check uses
\(q^{p-1}\equiv1\pmod p\) and \(q^{(p-1)/r}\not\equiv1\pmod p\) for every prime
divisor \(r\) of \(p-1\). These suffice since every proper divisor of \(p-1\)
divides \((p-1)/r\) for some such \(r\). For instance, at \(p=2791\) take \(q=53\).
Here \(2790=2\cdot3^2\cdot5\cdot31\), and the four exclusion residues in that
order are \(2790\), \(91\), \(2214\), and \(1000\). For more extensive tables
of least prime primitive roots, see Oliveira e Silva \cite{tos}.
\end{proof}

\begin{thebibliography}{9}
\bibitem{alp} L. Alp\"oge, \href{https://alpo.ge/remark.pdf}{announcement of the Lean proof of RH and Dirichlet GRH},
20 November 2026. \href{https://github.com/Deicyde/riemann-hypothesis}{Source and theorem statements}.
\bibitem{erdos} P. Erd\H{o}s, \emph{Some recent advances and current problems in number theory},
Lectures on Modern Mathematics, vol.~III (1965), 196--244, \href{https://www.renyi.hu/~p_erdos/1965-17.pdf\#page=39}{p.~233}.
\bibitem{prior} \texttt{lixiang90}, \href{https://github.com/TheJustinSunPrize/awards/pull/5615}{JSP-000818, pull request 5615},
The Justin Sun Prize, 20 November 2026.
\bibitem{mtt} K. McGown, E. Trevi\~no, and T. Trudgian,
\href{https://doi.org/10.7169/facm/2016.55.2.5}{\emph{Resolving Grosswald's conjecture on GRH}},
Funct. Approx. Comment. Math. \textbf{55} (2016), 215--225.
\bibitem{tos} T. Oliveira e Silva, \href{https://sweet.ua.pt/tos/p_roots.html}{\emph{Least prime primitive root of prime numbers}},
computational tables, University of Aveiro.
\bibitem{jan} N. Somani's GPT-Erdos project, \href{https://github.com/neelsomani/gpt-erdos/blob/9fead6532ee692f2a14d872175c52ddd19656698/data/solutions/985/candidate_solution.md}{candidate solution for problem 985}, 20 January 2026.
\bibitem{mit} A. V. Sutherland, \href{https://math.mit.edu/~drew/grh/}{\emph{Verification files for the Lean proof of the Riemann Hypothesis}},
MIT, 20 November 2026, 22:05 UTC. \href{https://math.mit.edu/~drew/grh/con-leche-compare.log}{Independent statement comparison}.
\end{thebibliography}
\end{document}
